\clearpage
\section{数值验证与模型评价}
\subsection{守恒、边界及解析解核验}
除场变量外，同时积分表面传质与热通量。离散质量审计采用
\begin{equation}
\sum_iV_i C_i(\t)+\int_0^{\t}\frac{\hmass}{\R}
 [\C(\R,s)-\Cinf(s)]\,ds
 =\sum_iV_i C_i(0).
\end{equation}
类似地，将$\rho\cp\R^2\sum_iV_i[T_i(\t)-T_i(0)]$与
$\int_0^{\t}\R\hconv[\Tinf(s)-\T(\R,s)]\,ds$比较。正式计算的相对质量收支残差为\texttt{\QOneMassError}，相对能量收支残差为\texttt{\QOneEnergyError}。这些结果证明离散收支闭合，不能代替对现实物理模型的实验验证。

另设置恒定环境温度阶跃及无浓度驱动的独立算例，将数值温度与具有对流边界的圆柱Bessel函数级数解对照；321节点算例的最大温度差小于$2\times10^{-4}$~K。测试还验证了均匀平衡态不变、热流方向正确、无驱动水分守恒以及Jacobian方向导数与直接扰动一致。

\subsection{空间与时间误差}
% 本文件由 scripts/make_tables.py 自动生成，请勿手工编辑
\begin{table}[htbp]
  \centering
  \caption{空间网格加密验证（七个指定时刻、21个输出位置的最大差）}
  \label{tab:validation}
  \begin{tabular}{ccc}
    \toprule
    节点数 & 相邻加密温度差/K & 相邻加密含水率差 \\
    \midrule
    81.0000 & -- & -- \\
    161.0000 & 3.18e-04 & 1.94e-04 \\
    321.0000 & 9.35e-05 & 4.86e-05 \\
    641.0000 & 2.11e-05 & 1.22e-05 \\
    \bottomrule
  \end{tabular}
\end{table}

表\ref{tab:validation}将嵌套网格的解插值到同一输出位置，在题目指定的七个时刻比较最大差。随网格加密，含水率误差差值约按四倍缩减，与二阶空间离散的预期一致。温度差还受到时间积分容差与空间插值的影响，不据此宣称每一个差值均严格满足二阶比例。

在641节点网格上将时间容差再收紧十倍，并将全程最大步长上限减半，指定时空点的温度最大变化为\texttt{\QOneTimeErrorT}~K，含水率最大变化为\texttt{\QOneTimeErrorC}~kg/kg。由此确认当前时间误差小于主要空间差值。表格保留四位小数仅为题目要求；接近舍入分界的末位仍可能因继续加密而变化。

\subsection{优点、局限与改进}
模型直接使用题给物性和实测环境序列，界面通量形式便于检验方向与守恒，表面加密提高了近表面水分场的分辨率。程序统一使用SI单位，物性、边界、求解器和导出相互分离，使结果可以复算并追溯来源。

主要局限是忽略端面、收缩以及蒸发冷却，同时采用有效浓度差闭合传质边界。若获得表面实测温度、气相湿度定义及吸附平衡数据，应先修订物理闭合并重新验证，再讨论预测精度，而不能仅通过加密计算网格消除这些结构性误差。
